\fsection{Punto de partida técnico}

El proyecto no parte de cero. Siemens proporciona, a través de su nota de aplicación "Robot Pick AI Pro Machine Vision Software for AI-based Universal Bin Picking", un proyecto base de TIA Portal y un archivo de instalación y programa para el robot UR. Estos ficheros, disponibles en Google Classroom, han sido verificados y probados en clase, y constituyen el punto de partida sobre el que se realizan las adaptaciones necesarias.
En concreto, los ficheros de partida son:

\begin{itemize}
	\item pickAi\_definitivo.URP: programa del robot UR desarrollado y validado en clase, que requiere ajuste de puntos de paso y posiciones de depósito.

	\item PickAI\_AppExample\_Installation.installation: archivo de instalación del robot UR.

	\item Proyecto TIA Portal: basado en el ejemplo de Siemens, con los tres FBs principales (1\_Sequence, 2\_PickAI y 3\_RobotCommunication) y el DB4 (MachineInterface) que centraliza los parámetros de configuración del sistema.
\end{itemize}
